\documentclass[10pt,a4paper]{article}
\usepackage[margin=1.8cm]{geometry}
\usepackage{booktabs}
\usepackage{array}
\usepackage{parskip}
\usepackage{enumitem}
\usepackage{titlesec}
\usepackage[T1]{fontenc}
\usepackage{fancyvrb}
\usepackage{xcolor}
\usepackage{mdframed}
\usepackage{hyperref}
\usepackage{multicol}

% ---- styling ----
\definecolor{codebg}{gray}{0.93}
\definecolor{notebg}{RGB}{255,248,220}
\definecolor{warnbg}{RGB}{255,235,235}

\newenvironment{codebox}{%
  \begin{mdframed}[backgroundcolor=codebg, linewidth=0pt,
    innerleftmargin=6pt, innerrightmargin=6pt,
    innertopmargin=4pt, innerbottommargin=4pt]%
  \small\ttfamily
}{%
  \end{mdframed}%
}

\newenvironment{notebox}{%
  \begin{mdframed}[backgroundcolor=notebg, linewidth=0.6pt,
    innerleftmargin=6pt, innerrightmargin=6pt,
    innertopmargin=3pt, innerbottommargin=3pt]%
  \small
}{%
  \end{mdframed}%
}

\newenvironment{warnbox}{%
  \begin{mdframed}[backgroundcolor=warnbg, linewidth=0.6pt,
    innerleftmargin=6pt, innerrightmargin=6pt,
    innertopmargin=3pt, innerbottommargin=3pt]%
  \small
}{%
  \end{mdframed}%
}

\newcommand{\cmd}[1]{\texttt{\small #1}}
\newcommand{\flag}[1]{\texttt{\small #1}}

\setlist[itemize]{noitemsep, topsep=2pt, leftmargin=*}
\setlist[enumerate]{noitemsep, topsep=2pt, leftmargin=*}

\titleformat{\section}{\large\bfseries}{}{0em}{}[\titlerule]
\titleformat{\subsection}{\normalsize\bfseries}{}{0em}{\textbullet\space}
\titlespacing{\section}{0pt}{10pt}{4pt}
\titlespacing{\subsection}{0pt}{6pt}{2pt}

\hypersetup{colorlinks=true, linkcolor=blue, urlcolor=blue}

\begin{document}

% ====================================================================
% CHEAT SHEET (1 page, tight margins)
% ====================================================================
\newgeometry{margin=0.55cm}
{\small                        % 9pt-ish body text for the cheat sheet
\setlength{\parskip}{1.5pt}
\setlength{\parindent}{0pt}
\titleformat{\section}{\normalfont\bfseries\normalsize}{}{0em}{}[\titlerule]
\titlespacing{\section}{0pt}{4pt}{2pt}

\begin{center}
  {\large\textbf{Git Cheat Sheet(from slide)}} \\[1pt]
  {\small From Basics to Branches}
\end{center}

\vspace{2pt}
\begin{multicols}{2}

%--- Section 1 ---
\section*{1. Setup \& Foundations}

\begin{tabular}{@{}p{3.2cm}p{4.5cm}@{}}
\toprule
\textbf{Command} & \textbf{What it actually does} \\
\midrule
\cmd{git init} & Creates \cmd{.git/} folder; stores entire project history. \\
\cmd{git status} & Shows branch, staged vs.\ modified files. \\
\cmd{.gitignore} & Prevents listed files from ever entering staging. \\
\bottomrule
\end{tabular}

%--- Section 2 ---
\section*{2. Staging \& Committing}

\textbf{Old workflow:}
\vspace{0.2em}
\begin{codebox}
    git add . \\ 
    git commit -m "did stuff" \\
    git push}
\end{codebox}

\textbf{Pro workflow:} Stage only files related to one feature, write a meaningful message.

\begin{tabular}{@{}p{3.4cm}p{2.7cm}p{1.2cm}@{}}
\toprule
\textbf{Command} & \textbf{Use when\ldots} & \textbf{New?} \\
\midrule
\cmd{git add -A} & Stage everything & \checkmark \\
\cmd{git commit -a} & Skip add for tracked files & $\times$ \\
\cmd{git commit -am "Msg"} & Skip add + inline msg & $\times$ \\
\cmd{git commit ----amend} & Fix last msg / forgot file. It will open a text editor and modify commited text in editor & N/A \\
\cmd{git commit ----amend -m "Edited msg"} & Change only the message & \\
\bottomrule
\end{tabular}

%--- Section 3 ---
\section*{3. Inspecting the Timeline}

\begin{itemize}
  \item \cmd{git diff} --- exact lines changed since last commit (unstaged).
  \item \cmd{git diff path/to/file} --- diff for one specific file.
  \item \cmd{git log} --- lists every commit hash + message (full timeline).
\end{itemize}

%--- Section 4 ---
\section*{4. Time Travel}

\begin{itemize}
  \item \cmd{git checkout path/to/file} --- discard changes; restore from last commit.
  \item \cmd{git checkout [hash] path/to/file} --- restore file from a specific commit.
  \item \cmd{git revert [hash]} --- new commit that undoes a past commit (safe).
\end{itemize}

\columnbreak

%--- Section 5 ---
\section*{5. Branches: Parallel Universes}

A branch is a pointer to a commit. \cmd{HEAD} points to your current branch.

\begin{tabular}{@{}p{3.8cm}p{3.8cm}@{}}
\toprule
\textbf{Command} & \textbf{Action} \\
\midrule
\cmd{git branch} & List all branches \\
\cmd{git branch [name]} & Create a new branch \\
\cmd{git checkout [name]} & Switch HEAD to that branch \\
\cmd{git checkout -b [name]} & Create \emph{and} switch (shortcut) \\
\cmd{git branch -d [name]} & Delete a branch \\
\bottomrule
\end{tabular}

\vspace{3pt}
\textbf{Detached HEAD:} Running \cmd{git checkout [hash]} points HEAD at a commit, not a branch --- you are floating in time!

\textbf{Fix:} \cmd{git checkout -b [new-branch-name]} to anchor it.

%--- Section 6 ---
\section*{6. Merging}

Git finds the least common ancestor to determine what changed.

\begin{enumerate}[noitemsep, topsep=1pt, leftmargin=*]
  \item \cmd{git checkout main} --- switch to target branch.
  \item \cmd{git merge feature-branch} --- pull in the changes.
  \item \cmd{git status} --- check for conflicts.
  \item Fix conflicting files, then \cmd{git add path/to/file}.
  \item \cmd{git commit} --- finalize the merge.
\end{enumerate}

%--- Section 7 ---
\section*{7. Remote Repositories (GitHub)}

\begin{itemize}
  \item \textbf{Git} = local time machine (runs in your terminal).
  \item \textbf{GitHub} = website for sharing \cmd{.git} repos.
  \item \cmd{git push origin main} --- upload local timeline so others can collaborate.
\end{itemize}

\end{multicols}
}% end \small group
\restoregeometry

% ====================================================================
% PRACTICAL WALKTHROUGH (original content from here)
% ====================================================================
\newpage

% Restore section/subsection formatting for the walkthrough
\titleformat{\section}{\large\bfseries}{}{0em}{}[\titlerule]
\titleformat{\subsection}{\normalsize\bfseries}{}{0em}{\textbullet\space}
\titlespacing{\section}{0pt}{10pt}{4pt}
\titlespacing{\subsection}{0pt}{6pt}{2pt}

% ---- Title ----
\begin{center}
  {\LARGE \textbf{Git \& GitHub: Every Command in a Real Project}}\\[4pt]
  {\normalsize A practical, narrative walkthrough --- one real folder, one real story}
\end{center}

\vspace{4pt}

\begin{notebox}
\textbf{Our project:} \texttt{F:\textbackslash Projects\textbackslash taskmanager} (Windows) ---
a small web app. We will build this project from zero to production, covering
\emph{every} command from the cheatsheet in the exact order you would naturally use them.
\end{notebox}

% ===================================================================
\section{Day 0 --- Starting Fresh: \texttt{git config} and \texttt{git init}}
% ===================================================================

You have just installed Git. Before touching any project, tell Git who you are.
This identity is stamped into every commit you ever make.

\begin{codebox}
git config ----global user.name "RiyadNabi"\\
git config ----global user.email "riyadunnabi@gmail.com"\\
git config ----global core.editor "code --wait"   \textcolor{gray}{\# use VS Code as editor\\}
\textcolor{gray}{\% git config ----global init.defaultBranch main \# otherwise Git's built-in default is master} \\
git config ----list   (or git config -l)                               \textcolor{gray}{\# verify everything is set}
\end{codebox}

My Setup:
\begin{codebox}
\tiny{\textit{PS C:\textbackslash Users\textbackslash riyad>} git config -l\\
diff.astextplain.textconv=astextplain\\
filter.lfs.clean=git-lfs clean -- \%f\\
filter.lfs.smudge=git-lfs smudge -- \%f\\
filter.lfs.process=git-lfs filter-process\\
filter.lfs.required=true\\
http.sslbackend=openssl\\
http.sslcainfo=C:/Program Files/Git/mingw64/etc/ssl/certs/ca-bundle.crt\\
core.autocrlf=true \quad core.fscache=true \quad core.symlinks=false\\
pull.rebase=false \quad credential.helper=manager\\
credential.https://dev.azure.com.usehttppath=true\\
init.defaultbranch=master\\
user.name=RiyadunNabi \\ user.email=riyadunnabi@gmail.com}
\end{codebox}

Now create and initialize the project folder:

\begin{codebox}
mkdir "F:\textbackslash Projects\textbackslash taskmanager"\\
% \hfill md    "F:\textbackslash Projects\textbackslash taskmanager"  \# Shortcut \\
cd    "F:\textbackslash Projects\textbackslash taskmanager"\\
git init -b main          \textcolor{gray}{\# start repo with 'main' as the default branch }
% cs\\
% (or, git init)
\end{codebox}

Or, more shortcut:

\begin{codebox}
md    "F:\textbackslash Projects\textbackslash taskmanager" \\
cd    "F:\textbackslash Projects\textbackslash taskmanager" \textcolor{gray}{\# (cd (Win\&Linux), chdir(identical on Win, different on Linux))} \\
git init          \textcolor{gray}{\# Uses whatever defaultBranch is set to}
\end{codebox}

Git creates a hidden \texttt{.git/} folder here. That folder \emph{is} Git's entire database ---
the full project history, branch pointers, everything. Your working files are untouched.

\begin{notebox}
\textbf{What just happened:} \cmd{git init} made \texttt{F:\textbackslash Projects\textbackslash taskmanager\textbackslash .git\textbackslash}.
\cmd{-b main} sets the initial branch name to \texttt{main} instead of the old default \texttt{master}.
\end{notebox}

% ===================================================================
\section{Day 0 --- First Files: \texttt{git status}, \texttt{.gitignore}, \texttt{git add}, \texttt{git commit}}
% ===================================================================

Create some starter files: \texttt{index.html}, \texttt{app.js}, \texttt{.env} (holds secrets),
and install Node dependencies (\texttt{node\_modules/}).

\begin{codebox}
git status          \textcolor{gray}{\# check what Git sees right now}\\
git status -s       \textcolor{gray}{\# compact version: same info, less noise}
\end{codebox}

Git now shows \texttt{index.html}, \texttt{app.js}, \texttt{.env}, \texttt{node\_modules/} as
\emph{untracked}. We never want to commit secrets or gigantic dependency folders.
Create \texttt{.gitignore}:

\begin{codebox}
\textcolor{gray}{\# .gitignore}\\
node\_modules/     \textcolor{gray}{\# ignore Node dependencies}\\
.env              \textcolor{gray}{\# ignore environment/secret files}\\
*.log             \textcolor{gray}{\# ignore all log files}\\
build/            \textcolor{gray}{\# ignore build output folder}\\
!.gitignore       \textcolor{gray}{\# but never ignore .gitignore itself}
\end{codebox}

Now stage and commit only the safe files:

\begin{codebox}
git add index.html app.js .gitignore\\
git status          \textcolor{gray}{\# should show 3 staged (green), .env and node\_modules gone}\\
git commit -m "Initial commit: add HTML, JS, and gitignore"
\end{codebox}

\begin{notebox}
\textbf{Staging area (index):} \cmd{git add} does not commit yet.
It moves file snapshots into the \emph{staging area}, a waiting room.
\cmd{git commit} then takes a permanent snapshot of whatever is in that room.
Think of it as: \emph{pack your suitcase} (add) then \emph{check it in} (commit).
\end{notebox}

% ===================================================================
\section{Day 1 --- Reviewing Work: \texttt{git log} and \texttt{git diff}}
% ===================================================================

Next morning you edit \texttt{app.js} and want to see exactly what changed before staging.

\begin{codebox}
git diff                    \textcolor{gray}{\# lines changed since last commit (unstaged)}\\
git diff app.js             \textcolor{gray}{\# same, but only for app.js}\\
git add app.js\\
git diff --staged           \textcolor{gray}{\# review what is NOW staged (ready to commit)}
\end{codebox}

Commit, then look at history:

\begin{codebox}
git commit -m "Add task creation logic"\\
git log                         \textcolor{gray}{\# full history with hashes, author, date, message}\\
git log --oneline               \textcolor{gray}{\# one line per commit --- easier to scan}\\
git log --graph --oneline --all \textcolor{gray}{\# ASCII branch/merge graph}\\
git log -p                      \textcolor{gray}{\# show actual diff introduced by each commit}
\end{codebox}

Each line in \cmd{git log --oneline} looks like: \texttt{a3f19c2 Add task creation logic}.
The 7-character prefix is the short SHA hash --- the permanent ID of that snapshot.

% ===================================================================
\section{Day 2 --- Branches: \texttt{git branch}, \texttt{git checkout}, \texttt{git switch}}
% ===================================================================

\subsection{Creating and switching branches}

You want to add a ``login'' feature without breaking the working \texttt{main} branch.

\begin{codebox}
git branch                      \textcolor{gray}{\# list all local branches (* = current)}\\
git branch feature/login        \textcolor{gray}{\# create the branch (HEAD still on main)}\\
git checkout feature/login      \textcolor{gray}{\# switch HEAD to that branch}\\
\textcolor{gray}{\# --- OR do both in one command: ---}\\
git checkout -b feature/login   \textcolor{gray}{\# create AND switch (most common)}\\
\textcolor{gray}{\# --- Modern equivalent: ---}\\
git switch -c feature/login     \textcolor{gray}{\# same as checkout -b, newer syntax}
\end{codebox}

\begin{notebox}
\textbf{What is HEAD?} HEAD is just a file inside \texttt{.git/} that contains the name of the
branch you are currently on. \cmd{git checkout feature/login} rewrites that file
and updates your working directory to match the branch's latest commit.
\end{notebox}

\subsection{Doing work on the branch}
\vspace{0.5em}
\begin{codebox}
\textcolor{gray}{\# ... edit login.html and auth.js ...}\\
git add -A                          \textcolor{gray}{\# stage all modified, deleted, and new files}\\
git commit -m "Add login page and auth module"\\
\\
\textcolor{gray}{\# Made a small typo in the message? Fix it before pushing:}\\
git commit --amend                  \textcolor{gray}{\# opens editor to rewrite last commit message}\\
\textcolor{gray}{\# Already have the right message, just forgot a file?}\\
git add forgotten.css\\
git commit --amend --no-edit        \textcolor{gray}{\# add file to last commit, keep message}
\end{codebox}

\subsection{Listing and deleting branches}
\vspace{0.5em}
\begin{codebox}
git branch -a           \textcolor{gray}{\# list local AND remote-tracking branches}\\
git branch -m login     \textcolor{gray}{\# rename current branch to 'login'}\\
git branch -d old-name  \textcolor{gray}{\# delete a fully merged branch (safe)}
\end{codebox}

% ===================================================================
\section{Day 3 --- The Detached HEAD Situation}
% ===================================================================

Curiosity strikes: you run \cmd{git log --oneline} and see an old commit hash \texttt{b2e7a1c}.
You do:

\begin{codebox}
git checkout b2e7a1c    \textcolor{gray}{\# browse old state of the project}
\end{codebox}

Git prints: \textit{``You are in 'detached HEAD' state.''}
HEAD now points directly at a commit, not at a branch.
Any new commits you make here will be \emph{orphaned} --- lost when you switch away.

\begin{warnbox}
\textbf{Fix immediately:} If you made changes you want to keep, anchor them to a real branch.\\
\texttt{git checkout -b snapshot-b2e7a1c}\\
If you just wanted to look and made nothing, simply: \texttt{git checkout main}
\end{warnbox}

% ===================================================================
\section{Day 4 --- Merging: \texttt{git merge} and Conflict Resolution}
% ===================================================================

The login feature is done. Merge it back into \texttt{main}:

\begin{codebox}
git checkout main                   \textcolor{gray}{\# always switch TO the receiving branch first}\\
git merge feature/login             \textcolor{gray}{\# pull the feature commits into main}\\
git merge --no-ff feature/login     \textcolor{gray}{\# force a merge commit even if fast-forward}\\
git status                          \textcolor{gray}{\# check for conflicts}
\end{codebox}

\textbf{If there are no conflicts:} Git creates a merge commit automatically. Done.

\textbf{If there are conflicts:} Git pauses and marks the conflicting files:

\begin{codebox}
\textcolor{gray}{\# Inside app.js you will see:}\\
<<<<<<< HEAD\\
const port = 3000;\\
=======\\
const port = 8080;\\
>>>>>>> feature/login
\end{codebox}

Resolve it manually --- delete the markers, keep the right code ---
then finalize:

\begin{codebox}
git add app.js              \textcolor{gray}{\# mark conflict as resolved}\\
git commit                  \textcolor{gray}{\# Git writes the merge commit message for you}\\
\textcolor{gray}{\# Changed your mind mid-merge?}\\
git merge --abort           \textcolor{gray}{\# abandon the merge, return to pre-merge state}
\end{codebox}

% ===================================================================
\section{Day 5 --- Comparing Branches: \texttt{git diff branch1..branch2}}
% ===================================================================

Before merging a second feature, preview what will change:

\begin{codebox}
git diff main..feature/dashboard    \textcolor{gray}{\# show all differences between two branches}
\end{codebox}

This is just \cmd{git diff} with two branch names. Extremely useful in code review.

% ===================================================================
\section{Day 6 --- Remote \& GitHub: \texttt{git remote}, \texttt{git push}, \texttt{git fetch}, \texttt{git pull}}
% ===================================================================

\subsection{Connecting to GitHub}

Create a blank repo on GitHub (do NOT initialize with a README). Then link it:

\begin{codebox}
git remote add origin https://github.com/riyad/taskmanager.git\\
git remote -v           \textcolor{gray}{\# verify: shows fetch and push URLs for 'origin'}
\end{codebox}

\subsection{Pushing for the first time}
\vspace{0.5em}
\begin{codebox}
git push -u origin main     \textcolor{gray}{\# -u sets 'origin/main' as upstream for future pulls}\\
git push --tags             \textcolor{gray}{\# push all local tags to remote}
\end{codebox}

After \cmd{-u}, future pushes on this branch only need \cmd{git push}.

\subsection{Cloning an existing repo}

On a different machine (or a teammate's):

\begin{codebox}
git clone https://github.com/riyad/taskmanager.git   \textcolor{gray}{\# full clone}\\
git clone --depth 1 https://github.com/riyad/taskmanager.git  \textcolor{gray}{\# shallow clone}\\
git clone -b feature/login https://github.com/riyad/taskmanager.git  \textcolor{gray}{\# specific branch}
\end{codebox}

\cmd{git clone} automatically sets \texttt{origin} to the URL. You get the full history.

\subsection{Fetching vs Pulling}

Your teammate pushed changes. How do you get them?

\begin{codebox}
git fetch                   \textcolor{gray}{\# download new commits; update origin/* refs; do NOT merge}\\
git fetch --prune           \textcolor{gray}{\# also remove stale origin/deleted-branch refs}\\
git log origin/main         \textcolor{gray}{\# inspect what they pushed before touching your work}\\
git diff main origin/main   \textcolor{gray}{\# compare your main to the remote one}\\
git pull                    \textcolor{gray}{\# fetch + merge in one step (can create a merge commit)}\\
git pull --rebase           \textcolor{gray}{\# fetch + rebase: keeps history linear (no merge commit)}
\end{codebox}

\begin{notebox}
\textbf{Fetch vs Pull:}
\cmd{git fetch} is safe --- it only updates the \texttt{origin/*} remote-tracking references.
\cmd{git pull} immediately applies those changes to your current branch.
Professional teams often prefer \cmd{fetch + inspect + merge} for more control.
\end{notebox}

\subsection{Force push (danger zone)}
\vspace{0.5em}
\begin{codebox}
git push --force    \textcolor{gray}{\# overwrite remote history --- use only on personal/feature branches}
\end{codebox}

\begin{warnbox}
\textbf{Never force-push to \texttt{main} or \texttt{develop}} on shared repos.
It rewrites history that teammates have already pulled, causing divergence for everyone.
\end{warnbox}

\subsection{Managing remotes}
\vspace{0.5em}
\begin{codebox}
git remote remove origin            \textcolor{gray}{\# remove the remote link (does not delete GitHub repo)}\\
git remote add upstream https://github.com/original/repo.git  \textcolor{gray}{\# common in open-source forks}
\end{codebox}

% ===================================================================
\section{Day 7 --- Rebase: \texttt{git rebase} and \texttt{git rebase -i}}
% ===================================================================

\subsection{Standard rebase (linearize history)}

Your \texttt{feature/dashboard} branch started from an old \texttt{main}.
You want it to look like it started from the \emph{current} tip of \texttt{main}:

\begin{codebox}
git checkout feature/dashboard\\
git rebase main             \textcolor{gray}{\# replay your commits on top of latest main}
\end{codebox}

Result: no merge commit, clean linear history. Great for pull requests.

\subsection{Interactive rebase (rewrite, squash, reorder)}

You made 5 messy commits while experimenting. Clean them up before merging:

\begin{codebox}
git rebase -i HEAD~5        \textcolor{gray}{\# open editor showing last 5 commits}
\end{codebox}

The editor shows: \vspace{0.2em}
\begin{codebox}
pick a1b2c3 Add dashboard skeleton\\
pick d4e5f6 wip\\
pick 7a8b9c more wip\\
pick 0d1e2f Fix typo\\
pick 3a4b5c Final dashboard logic
\end{codebox}

Change \texttt{pick} to \texttt{squash} (or \texttt{s}) on the wip commits to fold them in.
Save. Git re-applies and opens a final message editor.
Result: those 5 commits become 2 clean ones.

% ===================================================================
\section{Day 8 --- Undoing Things: \texttt{git reset} and \texttt{git revert}}
% ===================================================================

\subsection{Unstaging a file}

You accidentally staged \texttt{secret.env}:

\begin{codebox}
git reset secret.env        \textcolor{gray}{\# unstage, keep file on disk (default --mixed mode)}
\end{codebox}

\subsection{Undoing commits with reset}
\vspace{0.5em}
\begin{codebox}
git reset --soft HEAD~1     \textcolor{gray}{\# undo last commit, keep changes STAGED}\\
git reset HEAD~1            \textcolor{gray}{\# undo last commit, unstage changes (keep on disk)}\\
git reset --hard HEAD~1     \textcolor{gray}{\# undo last commit, DISCARD all changes (destructive)}
\end{codebox}

\begin{warnbox}
\cmd{--hard} is destructive. Changes are gone from disk. Never use it on pushed commits
unless you intend to force-push afterward.
\end{warnbox}

\subsection{Reverting a public commit}

You pushed a bad commit \texttt{c9f3a21} to \texttt{main}. You cannot rewrite history.
Use revert instead --- it creates a new commit that \emph{undoes} the old one:

\begin{codebox}
git revert c9f3a21      \textcolor{gray}{\# Git opens editor for the new "revert" commit message}\\
git push                \textcolor{gray}{\# safe to push: history is only added, never rewritten}
\end{codebox}

% ===================================================================
\section{Day 9 --- Saving Work Temporarily: \texttt{git stash}}
% ===================================================================

\textbf{Scenario:} You are halfway through \texttt{feature/dashboard} when your manager says
``fix a critical bug on \texttt{main} right now!'' Your changes are not ready to commit.

\begin{codebox}
git stash               \textcolor{gray}{\# save all staged + unstaged changes to a stash stack}\\
git stash list          \textcolor{gray}{\# see all stashes: stash@\{0\}, stash@\{1\}, ...}\\
git checkout main       \textcolor{gray}{\# now safe to switch branches (working tree is clean)}\\
\textcolor{gray}{\# ... fix the bug on main, commit, push ...}\\
git checkout feature/dashboard\\
git stash pop           \textcolor{gray}{\# restore stashed changes and remove them from the stack}\\
\textcolor{gray}{\# OR to apply but keep in stash:}\\
git stash apply stash@\{0\}\\
\textcolor{gray}{\# Remove a specific stash:}\\
git stash drop          \textcolor{gray}{\# removes top stash (stash@\{0\})}
\end{codebox}

% ===================================================================
\section{Day 10 --- Restoring Files: \texttt{git checkout -{}- file}}
% ===================================================================

You edited \texttt{app.js} but it's a mess and you want to throw away all changes since the last commit:

\begin{codebox}
git checkout -- app.js      \textcolor{gray}{\# restore app.js from last commit (discard working changes)}
\end{codebox}

\begin{warnbox}
There is \textbf{no undo} for this. The working-directory changes are gone.
\end{warnbox}

% ===================================================================
\section{Day 11 --- Tagging Releases: \texttt{git tag}}
% ===================================================================

Version 1.0 is ready. Tag the release commit:

\begin{codebox}
git tag v1.0                            \textcolor{gray}{\# lightweight tag (just a pointer)}\\
git tag -a v1.0 -m "First production release"   \textcolor{gray}{\# annotated tag (stores author + date)}\\
git tag                                 \textcolor{gray}{\# list all tags}\\
git push --tags                         \textcolor{gray}{\# push all local tags to remote}\\
git tag -d v1.0                         \textcolor{gray}{\# delete a tag locally}
\end{codebox}

Tags differ from branches: they never move when new commits are added.
They are permanent markers for releases.

% ===================================================================
\section{Unconventional Scenario A --- Moving Uncommitted Work to Another Branch}
% ===================================================================

You started coding on \texttt{main} by mistake. The changes are not committed yet.

\begin{codebox}
git stash                   \textcolor{gray}{\# 1. save the accidental work}\\
git checkout -b dev         \textcolor{gray}{\# 2. create the correct branch}\\
git stash pop               \textcolor{gray}{\# 3. reapply the changes here}\\
git add .\\
git commit -m "Work correctly started on dev"
\end{codebox}

% ===================================================================
\section{Unconventional Scenario B --- Removing a Secret File That Was Committed}
% ===================================================================

You committed \texttt{.env} (with real passwords) and pushed it. Act fast:

\begin{codebox}
git rm --cached .env                    \textcolor{gray}{\# untrack the file, leave it on disk}\\
echo ".env" >> .gitignore               \textcolor{gray}{\# make sure it never comes back}\\
git commit -m "Remove .env from tracking"\\
git push
\end{codebox}

\begin{notebox}
The file is gone from \emph{future} commits, but it is still in old commits in history.
For a real secret leak, rotate the credentials AND use a tool like
\texttt{git filter-repo} to scrub history entirely.
\end{notebox}

% ===================================================================
\section{Unconventional Scenario C --- Recovering a Lost Commit (\texttt{git reflog})}
% ===================================================================

You ran \cmd{git reset --hard} and realized you needed that commit.
Git almost never truly deletes data. The reflog is your safety net:

\begin{codebox}
git reflog show             \textcolor{gray}{\# every HEAD movement Git has recorded}\\
\textcolor{gray}{\# Output example:}\\
\textcolor{gray}{\# a3f19c2 HEAD@\{0\}: reset: moving to HEAD~1}\\
\textcolor{gray}{\# 9b2c41a HEAD@\{1\}: commit: Add payment module  <-- this is what you lost}\\
\\
git checkout -b recover 9b2c41a     \textcolor{gray}{\# restore lost commit on a new branch}
\end{codebox}

% ===================================================================
\section{Unconventional Scenario D --- Cherry-Picking One Commit}
% ===================================================================

The \texttt{feature/payments} branch has one specific hotfix commit (\texttt{ff1a390})
that \texttt{main} urgently needs, but the full branch is not ready to merge.

\begin{codebox}
git checkout main\\
git cherry-pick ff1a390     \textcolor{gray}{\# copy ONLY that commit onto main}
\end{codebox}

This does not merge the whole branch --- just replays the patch of that one commit.

% ===================================================================
\section{Unconventional Scenario E --- Finding a Bug with \texttt{git bisect}}
% ===================================================================

The app worked at tag \texttt{v1.0} but is broken now. Which commit introduced the bug?

\begin{codebox}
git bisect start\\
git bisect bad                  \textcolor{gray}{\# current HEAD is broken}\\
git bisect good v1.0            \textcolor{gray}{\# v1.0 was working}\\
\textcolor{gray}{\# Git checks out the midpoint commit. Test the app. Then:}\\
git bisect good                 \textcolor{gray}{\# if the bug is absent at this midpoint}\\
git bisect bad                  \textcolor{gray}{\# if the bug is present}\\
\textcolor{gray}{\# Repeat ~log2(N) times. Git will print: "X is the first bad commit"}\\
git bisect reset                \textcolor{gray}{\# return HEAD to original position}
\end{codebox}

If your test is scriptable, you can automate the entire bisect:

\begin{codebox}
git bisect run ./test.sh        \textcolor{gray}{\# script must exit 0 for good, non-zero for bad}
\end{codebox}

% ===================================================================
\section{Unconventional Scenario F --- Interactively Staging Hunks: \texttt{git add -p}}
% ===================================================================

You made two unrelated edits in \texttt{app.js}: a bug fix and a new feature.
You want two separate commits for clarity.

\begin{codebox}
git add -p app.js       \textcolor{gray}{\# Git shows each changed "hunk" and asks: stage it? (y/n/s/q)}
\end{codebox}

Press \texttt{y} to stage a hunk, \texttt{n} to skip, \texttt{s} to split a hunk further.
Commit the staged portion, then \cmd{git add -p} again for the rest.

% ===================================================================
\section{Unconventional Scenario G --- \texttt{git diff} Between Branches and Commits}
% ===================================================================

\begin{codebox}
git diff main..feature/login            \textcolor{gray}{\# what does the feature add on top of main?}\\
git diff a3f19c2..9b2c41a              \textcolor{gray}{\# diff between any two commit hashes}\\
git diff HEAD~3                         \textcolor{gray}{\# compare working tree to 3 commits ago}
\end{codebox}

% ===================================================================
\section{Quick Reference: Commands Grouped by Phase}
% ===================================================================

\begin{tabular}{@{}p{4.5cm}p{10cm}@{}}
\toprule
\textbf{Phase} & \textbf{Commands used} \\
\midrule
Setup & \cmd{git config --global}, \cmd{git init -b}, \cmd{git clone}, \cmd{git clone --depth 1}, \cmd{git clone -b} \\[2pt]
Inspect & \cmd{git status}, \cmd{git status -s}, \cmd{git log}, \cmd{git log --oneline}, \cmd{git log --graph --oneline --all}, \cmd{git log -p}, \cmd{git diff}, \cmd{git diff --staged}, \cmd{git diff branch1..branch2} \\[2pt]
Stage \& Commit & \cmd{git add [file]}, \cmd{git add .}, \cmd{git add -A}, \cmd{git add -p}, \cmd{git commit -m}, \cmd{git commit -a -m}, \cmd{git commit --amend}, \cmd{git commit --amend --no-edit} \\[2pt]
Branches & \cmd{git branch}, \cmd{git branch -a}, \cmd{git branch [name]}, \cmd{git branch -d}, \cmd{git branch -m}, \cmd{git checkout [branch]}, \cmd{git checkout -b}, \cmd{git checkout --}, \cmd{git switch}, \cmd{git switch -c} \\[2pt]
Merge \& Rebase & \cmd{git merge}, \cmd{git merge --no-ff}, \cmd{git merge --abort}, \cmd{git rebase}, \cmd{git rebase -i}, \cmd{git cherry-pick} \\[2pt]
Remote & \cmd{git remote add}, \cmd{git remote -v}, \cmd{git remote remove}, \cmd{git fetch}, \cmd{git fetch --prune}, \cmd{git pull}, \cmd{git pull --rebase}, \cmd{git push}, \cmd{git push -u}, \cmd{git push --force}, \cmd{git push --tags} \\[2pt]
Undo & \cmd{git reset [file]}, \cmd{git reset --soft HEAD\textasciitilde1}, \cmd{git reset --hard}, \cmd{git revert}, \cmd{git stash}, \cmd{git stash pop}, \cmd{git stash list}, \cmd{git stash apply}, \cmd{git stash drop}, \cmd{git rm --cached}, \cmd{git reflog} \\[2pt]
Tags & \cmd{git tag}, \cmd{git tag -a}, \cmd{git tag -d}, \cmd{git push --tags} \\[2pt]
Debug & \cmd{git bisect start/bad/good/reset/run}, \cmd{git add -p} \\[2pt]
Config & \cmd{git config --global user.name/email/core.editor}, \cmd{git config --list} \\
\bottomrule
\end{tabular}

\vspace{6pt}
\begin{notebox}
\textbf{Mental model summary:}
Your project lives in three places simultaneously:
(1) \textbf{Working directory} --- what you see in your file manager.
(2) \textbf{Staging area (index)} --- what \cmd{git add} has prepared for the next commit.
(3) \textbf{Repository (\texttt{.git/})} --- permanent snapshots recorded by \cmd{git commit}.
Remote is a fourth location (GitHub) synced by \cmd{push/pull/fetch}.
Every Git command is fundamentally about moving data between these four places.
\end{notebox}

\end{document}